\section{Контрольная по занятиям №1--3}\label{solutions:mock13}
Максимум --- 100 баллов: 60 за В1--В5 и 40 за В6--В8. Баллы за метод начисляются при правильном ходе решения, в том числе при отдельной арифметической ошибке. В каждом подпункте учитывайте указанный максимум.

\subsection*{В1. Переводы и арифметика}\label{sol:v1}
\hyperref[task:v1]{К условию.}
$205=128+64+8+4+1$, поэтому двоичная запись $11001101_2$. Группы \texttt{1100 1101} дают $CD_{16}$. Проверка: $12\cdot16+13=205$.

Для восьмеричного сложения начинаем справа: $7+6=13_{10}=15_8$, пишем 5, переносим 1. Следующий разряд: $5+2+1=8_{10}=10_8$, пишем 0, переносим 1 в новый старший разряд. Ответ $105_8$.
Проверка: $57_8=47$, $26_8=22$, $47+22=69=105_8$.

\textbf{Баллы:} подпункт 1 --- 2 за способ перевода, 2 за двоичную и 2 за hex-запись; подпункт 2 --- 2 за разрядные действия с переносами, 1 за ответ и 1 за десятичную проверку. Всего 10.

\subsection*{В2. Диапазоны и четыре кода}\label{sol:v2}
\hyperref[task:v2]{К условию.}
Для 7-битного дополнительного кода диапазон равен
\[
[-2^6,2^6-1]=[-64,63].
\]
Количество значений $2^7=128$.

Для кодирования $-27$ переходим к ширине 8 битов. Положительное $27=16+8+2+1$ имеет запись \texttt{00011011}.
\begin{itemize}
    \item Прямой код: знак 1 и семибитный модуль \texttt{0011011}. Итог: \texttt{10011011}, $9B_{16}$.
    \item Обратный код: инвертируем восемь битов положительной записи. Получаем \texttt{11100100}, $E4_{16}$.
    \item Дополнительный код: прибавляем 1 к обратному. Получаем \texttt{11100101}, $E5_{16}$. Проверка: $256-27=229$.
    \item Со смещением: $-27+128=101$, поэтому код \texttt{01100101}, $65_{16}$. Проверка: $101-128=-27$.
\end{itemize}
\textbf{Баллы:} за диапазон --- по 0.5 за каждую границу, за количество значений --- 1. За каждый из четырёх кодов --- 1 за получение и 1 за согласованные двоичную и hex-записи. Всего 10.

\subsection*{В3. Результат и флаги}\label{sol:v3}
\hyperref[task:v3]{К условию.}
\textbf{Сложение.} $85=55_{16}$, $60=3C_{16}$. Сумма:
\[
55_{16}+3C_{16}=91_{16}=\texttt{10010001}.
\]
Сохранённое знаковое значение $145-256=-111$. CF=0, поскольку беззнаковая сумма $145<256$. OF=1, поскольку точная знаковая сумма $145>127$. Результат ненулевой, старший бит 1: ZF=0, SF=1.

\textbf{Вычитание.} $40=28_{16}$, $75=4B_{16}$. Дополнительный код $-75$ получается инверсией и прибавлением единицы:
\[
\texttt{01001011}\ \longrightarrow\ \texttt{10110100}
\ \longrightarrow\ \texttt{10110101}.
\]
Сумма $28_{16}+B5_{16}=DD_{16}=\texttt{11011101}$. Знаковое значение $221-256=-35$ совпадает с точной разностью.

Для SUB заём нужен, поскольку $40<75$: CF=1. Разность $-35$ представима: OF=0. ZF=0, SF=1. Вспомогательное сложение $40+181=221$ выполняется без переноса; CF исходного вычитания определяется по займу.

\textbf{Баллы за каждую операцию:} 1 за код результата, 1 за знаковое значение, по 2 за CF и OF с обоснованиями, по 0.5 за ZF и SF. За получение разности через дополнительный код --- ещё 1 балл. Всего $7+8=15$.

\subsection*{В4. Операции над битами}\label{sol:v4}
\hyperref[task:v4]{К условию.}
Исходное слово $B3_{16}=\texttt{10110011}$. Все действия выполняются над ним независимо.
\begin{enumerate}
    \item $5C_{16}=\texttt{01011100}$. Общая единица находится только в бите 4, поэтому AND даёт \texttt{00010000}, $10_{16}$.
    \item Для сброса бита 5 нужна маска $DF_{16}=\texttt{11011111}$. AND с этой маской даёт \texttt{10010011}, $93_{16}$.
    \item Бит 2 устанавливаем через OR с маской $04_{16}$. Её биты: \texttt{00000100}. Ответ: \texttt{10110111}, $B7_{16}$.
    \item Извлекаем три бита начиная с номера 2. Логический сдвиг вправо на 2 даёт $2C_{16}=\texttt{00101100}$. AND с маской $07_{16}$ оставляет \texttt{00000100}, то есть $04_{16}$, число 4.
\end{enumerate}
\textbf{Баллы:} в подпункте 1 --- 1 за побитовое вычисление и 1 за ответ; в подпунктах 2 и 3 --- по 1 за операцию с маской и по 1 за ответ; в подпункте 4 --- по 1 за сдвиг, маску с AND, код результата и десятичное значение. Всего 10.

\subsection*{В5. Сдвиги и расширение}\label{sol:v5}
\hyperref[task:v5]{К условию.}
Исходное $D5_{16}=\texttt{11010101}$ имеет unsigned-значение 213 и знаковое значение $213-256=-43$.
\begin{enumerate}
    \item Логический сдвиг даёт $35_{16}=53$, биты \texttt{00110101}. Проверка: $\lfloor213/4\rfloor=53$.
    \item При арифметическом сдвиге слева дописываем две единицы. Результат: \texttt{11110101}, $F5_{16}$. Значение $245-256=-11$.
    \item Циклический сдвиг влево на 1 переносит старшую единицу в младший разряд: \texttt{10101011}, $AB_{16}$.
    \item Знаковое расширение повторяет старший бит:
    \[
    \texttt{11111111 11010101}=FFD5_{16}.
    \]
    \item Точное частное $-43/4=-10.75$. При округлении к нулю получаем $-10$, при округлении вниз --- $-11$. Арифметический сдвиг реализует второй способ округления.
\end{enumerate}
\textbf{Баллы:} в подпунктах 1 и 2 --- по 1 за биты, hex и значение; в подпунктах 3 и 4 --- по 1 за правило, биты и hex; в подпункте 5 --- по 1 за исходное знаковое значение, частное и объяснение различия. Всего 15.

\subsection*{В6. Представление float32}\label{sol:v6}
\hyperref[task:v6]{К условию.}
\textbf{Кодирование.} Целая часть $13=1101_2$; дробная $0.375=0.25+0.125=0.011_2$. Нормализуем:
\[
-13.375=-1101.011_2=-1.101011_2\cdot2^3.
\]
Знак $S=1$. Истинный порядок $E=3$, хранимый порядок $e=3+127=130=10000010_2$. Дробное поле содержит биты после ведущей единицы; дополняем их до 23 битов:
\begin{center}
\texttt{10101100000000000000000}.
\end{center}
Собранное слово и его группировка:
\begin{center}
\texttt{1 10000010 10101100000000000000000}\\
\texttt{1100 0001 0101 0110 0000 0000 0000 0000}.
\end{center}
Ответ \texttt{C1560000}. Двоичная дробь конечна и помещается в формат, поэтому округление значения не требуется.

\textbf{Декодирование.} Для \texttt{40B80000} поля имеют вид:
\begin{center}
\texttt{0 10000001 01110000000000000000000}.
\end{center}
Знак $S=0$, порядок $e=129$. Число нормализованное, истинный порядок $E=129-127=2$. Значащая часть со скрытой единицей:
\[
1.0111_2=1+2^{-2}+2^{-3}+2^{-4}=\frac{23}{16}.
\]
Значение $\frac{23}{16}\cdot2^2=\frac{23}{4}=5.75$.

\textbf{Баллы:} кодирование --- 2 за перевод и нормализацию, 1 за знак, 2 за истинный и хранимый порядки с битовой записью, 2 за дробное поле и сборку слова, 2 за hex. Декодирование --- по 1 за каждое из трёх полей, 1 за истинный порядок, 1 за значащую часть и 1 за значение. Всего $9+6=15$.

\subsection*{В7. Особые и денормализованные значения}\label{sol:v7}
\hyperref[task:v7]{К условию.}
Обозначим через $f$ целочисленное значение 23-битного дробного поля.
\begin{enumerate}
    \item \texttt{80000000}: $S=1$, $e=0$, $f=0$. Это отрицательный ноль $-0$.
    \item \texttt{7F800000}: $S=0$, $e=255$, $f=0$. Это положительная бесконечность $+\infty$.
    \item \texttt{7FC00001}: $S=0$, $e=255$, $f=400001_{16}\ne0$. Это NaN. Бит знака равен 0; числового значения у NaN нет.
    \item \texttt{00400000}: $S=0$, $e=0$, $f=400000_{16}=2^{22}\ne0$. Это положительное денормализованное число. В значащей части перед точкой стоит 0, а масштаб задаётся степенью $2^{-126}$:
    \[
    x=0.1_2\cdot2^{-126}
      =\frac{2^{22}}{2^{23}}\cdot2^{-126}=2^{-127}.
    \]
\end{enumerate}
\textbf{Баллы:} за каждый из первых трёх подпунктов --- 1 за поля и 1 за класс. За четвёртый --- 1 за поля, 1 за класс, 1 за формулу с нулевым скрытым битом и масштабом $2^{-126}$, 1 за точное значение. Всего 10.

\subsection*{В8. Округление и неассоциативность}\label{sol:v8}
\hyperref[task:v8]{К условию.}
У нормализованного \texttt{float32} 23 бита дробной части. Число 1 имеет истинный порядок 0. При увеличении дробного поля на единицу следующая запись даёт $1+2^{-23}$. Поэтому
\[
\varepsilon=2^{-23},\qquad h=2^{-24}=\frac{\varepsilon}{2}.
\]

\textbf{Левая расстановка скобок.} Точная сумма $1+h$ лежит посередине между соседними числами 1 и $1+\varepsilon$. У числа 1 младший бит значащей части равен 0, у следующего числа --- 1. Правило округления к чётному выбирает 1. При втором сложении возникает та же ситуация:
\[
\operatorname{fl}(1+h)=1,\qquad
\operatorname{fl}(\operatorname{fl}(1+h)+h)=1.
\]
Здесь $\operatorname{fl}$ обозначает округление до \texttt{float32}. Код результата \texttt{3F800000}.

\textbf{Правая расстановка скобок.} Сначала складываем равные малые числа:
\[
h+h=2^{-23}=\varepsilon.
\]
Эта сумма точно представима в \texttt{float32}. Следующее сложение также даёт представимое число:
\[
\operatorname{fl}(1+\operatorname{fl}(h+h))
=1+\varepsilon=1+2^{-23}.
\]
Его код \texttt{3F800001}. Результаты двух вычислений различаются на $2^{-23}$ --- один шаг между соседними числами около 1.

В левом выражении малое слагаемое дважды поглощается при округлении. Различные ответы при перестановке скобок показывают потерю ассоциативности сложения. Точная сумма обоих выражений равна $1+2h=1+\varepsilon$; различие возникло при округлении промежуточных результатов.

\textbf{Баллы:} подпункт 1 --- 1 за связь с 23 битами дробного поля, 1 за следующее число и 1 за $\varepsilon$; подпункт 2 --- по 2 за каждый шаг с округлением; подпункт 3 --- по 2 за каждый точный шаг; подпункт 4 --- 1 за положение посередине, 1 за выбор чётного по младшему биту, по 1 за объяснение поглощения и потери ассоциативности. Всего 15.

\section{Повторение по занятиям №1--3}\label{review:mock13}
Материал для повторения приведён в основном пособии \href{adaptation_course.pdf}{«Адаптационный курс»}. В таблице указаны соответствующие занятия и задания практикума.

\begin{center}
\small
\begin{tabular}{>{\raggedright\arraybackslash}p{0.31\textwidth}|>{\raggedright\arraybackslash}p{0.35\textwidth}|>{\raggedright\arraybackslash}p{0.23\textwidth}}
Ошибка & Материал курса & Задания практикума\\\hline
Ошибки перевода и арифметики в заданном основании & \href{adaptation_course.pdf}{Лекция и семинар №1}: системы счисления и арифметика & \hyperref[task:v1]{В1}\\[4pt]
Неверный диапазон или код целого числа & \href{adaptation_course.pdf}{Лекция и семинар №2}: диапазоны и коды & \hyperref[task:v2]{В2}\\[4pt]
Смешение CF и OF, ошибка в определении займа & \href{adaptation_course.pdf}{Лекция и семинар №2}: арифметика и флаги & \hyperref[task:v3]{В3}\\[4pt]
Маска меняет лишние биты & \href{adaptation_course.pdf}{Лекция и семинар №2}: побитовые операции & \hyperref[task:v4]{В4}\\[4pt]
Неверное заполнение битов при сдвиге или расширении & \href{adaptation_course.pdf}{Лекция и семинар №2}: сдвиги и знаковое расширение & \hyperref[task:v5]{В5}\\[4pt]
Ошибка в нормализации, порядке или дробном поле & \href{adaptation_course.pdf}{Лекция и семинар №3}: представление IEEE 754 & \hyperref[task:v6]{В6}\\[4pt]
Смешение особых и денормализованных значений & \href{adaptation_course.pdf}{Лекция и семинар №3}: классы значений \texttt{float32} & \hyperref[task:v7]{В7}\\[4pt]
Ошибка в машинном эпсилоне или округлении промежуточных результатов & \href{adaptation_course.pdf}{Лекция и семинар №3}: точность, округление и поглощение & \hyperref[task:v8]{В8}
\end{tabular}
\end{center}
